\chapter{La reforma y su reversión}
\label{ch:dollarreform}

\section{Reforma bajo presión}

Entre julio de 1993 y finales de 1995 el gobierno cubano hizo más por liberalizar
su economía que en los treinta y cuatro años anteriores juntos. Lo hizo a
regañadientes, describió las medidas en todo momento como concesiones temporales
forzadas por las circunstancias, y retiró tantas como pudo en cuanto estuvo en
condiciones de hacerlo. Esa secuencia es el asunto de este capítulo, y es el
ejemplo más nítido del patrón que este libro ha venido armando.

Las medidas, en orden:

\emph{Julio de 1993.} Se despenalizó la tenencia de divisa. Hasta ese mes, un
cubano al que se le encontrara un dólar podía ser procesado; desde ese mes los
cubanos podían tener, gastar y recibir dólares legalmente. Fue la decisión
bisagra, porque legalizó las remesas ---y las remesas de la diáspora que treinta
años de política habían denunciado como gusanos se convirtieron, en una década, en
una de las mayores fuentes de divisas del país---.

\emph{Septiembre de 1993.} Se relegalizó el trabajo por cuenta propia en 117
categorías, luego ampliadas a unas 157. La lista era una lista: la actividad se
permitía por enumeración, no por derecho, y quedaron excluidas las profesiones de
formación universitaria ---un médico no podía ejercer privadamente, un ingeniero
no podía asesorar, pero cualquiera de los dos podía alquilar un cuarto o manejar
un taxi---. Esta decisión de diseño, más que ninguna otra, determinó en qué se
convirtió el sector privado cubano.

\emph{Septiembre de 1993.} Las granjas estatales, que tenían la mayor parte de la
tierra agrícola, se fragmentaron en Unidades Básicas de Producción Cooperativa,
las UBPC, con la tierra en usufructo y un derecho nominal a organizar su propio
trabajo. Cerca del 40 por ciento de la tierra agrícola cambió de manos sobre el
papel.\unv{} En la práctica el Estado conservó el monopolio de acopio, el
suministro de insumos y el precio, de modo que lo que se descentralizó fue la
pérdida y no la decisión. Las UBPC han rendido menos que el sector privado y que
las cooperativas más antiguas de forma continuada desde entonces.

\emph{Octubre de 1994.} Se reabrieron los mercados agropecuarios a precios libres
para la producción por encima de la cuota estatal obligatoria ---la misma medida
que se había introducido en 1980 y abolido en 1986---. En semanas apareció comida
en La Habana, a precios que la mayoría de la gente no podía pagar, que fue también
lo que había ocurrido en 1980.

\emph{1994--95.} El ajuste fiscal y monetario, que es la parte de este período
que menos atención recibe y la que más merece. El déficit presupuestario, que
había llegado a algo así como el 30 por ciento del producto en 1993, se recortó a
alrededor del 7 por ciento en 1994 y al 2 por ciento para 1996, mediante subidas
de precios, recortes de subsidios, nuevos impuestos bajo una ley tributaria de
1994, y la retirada de circulante por medios confiscatorios. En 1995 se abrieron
casas de cambio a tasa de mercado. La tasa informal, que había llegado a 120 o 150
pesos por dólar a mediados de 1994, volvió a unos 20 o 25 para 1996 y se mantuvo
en esa banda los veinticinco años siguientes \citep{mesalago2005}.\unv{} Una
moneda que había sido destruida se estabilizó en dos años sin asistencia externa.
El capítulo~\ref{ch:money} toma esto en serio como precedente, porque el mismo
Estado hizo lo contrario en 2021.

\emph{Septiembre de 1995.} La Ley 77 abrió la inversión extranjera, permitiendo
en principio hasta la propiedad totalmente extranjera y creando la forma de
empresa mixta que ha usado desde entonces la mayor parte del capital foráneo en
Cuba. Conservó el rasgo que más ha hecho por disuadir a los inversionistas: el
socio extranjero no contrata a sus propios trabajadores, sino que los toma de una
agencia empleadora estatal, le paga a la agencia en divisa, y la agencia le paga
al trabajador en pesos, quedándose con la diferencia. A una tasa oficial de uno
por uno, eso significaba que un trabajador por el que su empleador pagaba al
Estado varios cientos de dólares al mes recibía el equivalente en pesos de unos
veinte. El capítulo~\ref{ch:capital} sostiene que ese único arreglo le ha costado
a Cuba más inversión extranjera que cualquier estatuto norteamericano.

Las medidas funcionaron. 1994 fue el fondo; el crecimiento se reanudó en 1995 y
corrió durante la segunda mitad de la década. La economía que emergió, sin
embargo, no era la que había entrado.

\section{La economía dual}

Cuba salió del Período Especial con dos monedas, dos niveles de precios, dos
mercados de trabajo y dos poblaciones.

El peso convertible, el CUC, se introdujo en 1994 y se fijó a uno por dólar. Junto
a él siguió el peso corriente, a veinticuatro o veinticinco por CUC para los
hogares y ---esta es la parte que importaba--- a uno por uno para la contabilidad
de las empresas estatales. Cuba mantuvo así, durante veintiséis años, una
discrepancia de veinticinco veces entre la tasa que enfrentaban sus hogares y la
que sus propias empresas usaban para valorar importaciones, exportaciones, costos
y utilidades. Toda la contabilidad empresarial del país era ficción. Ningún
gerente cubano podía saber si su operación creaba o destruía valor, y el
ministerio por encima tampoco. Esto no es un tecnicismo; es la razón de que la
industria cubana no pudiera reformarse ni siquiera por quienes querían
reformarla, porque nadie podía identificar qué partes valía la pena conservar.

La consecuencia social fue la inversión de la pirámide. Un cirujano ganaba un
salario estatal en pesos. Un botones ganaba propinas en divisa. Una profesora de
literatura manejaba un taxi por la tarde; una química investigadora alquilaba su
cuarto libre. La valoración que la sociedad había pasado treinta años enseñando
---que la educación era la vía al reconocimiento--- quedaba contradicha a diario
y en público por el tipo de cambio. El coeficiente de Gini de Cuba, que había
estado cerca de 0,24 a mediados de los ochenta y era de los más bajos medidos en
cualquier parte, andaba por 0,40 al final de los noventa
\citep{mesalago2005}.\unv{}

Esa desigualdad tenía color. La emigración de 1959--62 había sido abrumadoramente
blanca; las remesas que volvían treinta años después fluían por tanto
abrumadoramente hacia hogares blancos, en proporciones que las encuestas
disponibles sitúan en cuatro quintas partes o más.\unv{} El sector turístico, la
otra vía a la divisa, contrataba al personal de cara al público por ``buena
presencia'', que en la práctica significaba lo que siempre ha significado. De modo
que el Estado que había abolido la barrera de color en 1959 y luego declarado
cerrada la cuestión racial en 1962 ---volviéndola así inmedible durante treinta
años--- entró en su apertura de mercado con el acceso a la divisa distribuido
exactamente a lo largo de la línea que había dejado de mirar. El
capítulo~\ref{ch:premises} trata esto como restricción dura de cualquier diseño de
transición, y no como nota histórica.

El turismo fue la cara visible. Las llegadas pasaron de unas 340.000 en 1990 a
unas 745.000 en 1995 y a cerca de 1,8 millones en 2000.\unv{} La planta la
construyeron empresas mixtas con cadenas españolas, canadienses y jamaicanas,
sobre todo como resorts de playa todo incluido en Varadero y los cayos del norte,
físicamente separados del país que los rodeaba ---y, hasta 2008, también
legalmente separados, puesto que a los cubanos no se les permitía entrar en los
hoteles turísticos---. Un Estado que había fundado su reclamo sobre la abolición
de una Cuba para extranjeros había reconstruido una, y esta vez el gobierno era
quien vigilaba la puerta.

\section{La reversión}

Para 1996 la emergencia había pasado, y con ella cambió la evaluación que la
dirigencia hacía de las reformas.

La señal llegó en marzo de 1996, cuando Raúl Castro presentó al Comité Central un
informe que atacaba al Centro de Estudios sobre América, un instituto de
investigación del partido cuyos economistas venían analizando las reformas y su
posible extensión. El instituto se desmanteló y sus investigadores se dispersaron.
El mensaje a las ciencias sociales cubanas se recibió y nunca ha necesitado
repetirse.

Lo que siguió fue una década de reversión callada. Las licencias de trabajo por
cuenta propia, que habían llegado a unas 209.000 a comienzos de 1996, se
estrujaron con impuestos, inspecciones y retirada de categorías, cayendo a unas
150.000 en 2003 y menos después \citep{ritter2015}.\unv{} Las nuevas licencias en
la mayoría de las categorías se suspendieron del todo en 2003--04. En 2003 se
obligó a las empresas a mantener y operar su divisa a través de una única cuenta
central, poniendo fin al acceso descentralizado que había hecho funcionar al
sector de empresas mixtas. En noviembre de 2004 el dólar estadounidense se retiró
por completo de la circulación y se impuso un recargo del 10 por ciento a la
conversión de dólares en CUC, medida presentada como respuesta a la presión
financiera norteamericana y que además gravaba cada remesa.

Nada de esto fue un regreso a 1968; los mercados siguieron abiertos, el turismo
creció, las remesas siguieron llegando. Pero la dirección era inconfundible, y su
lógica era coherente: las reformas de 1993--95 habían sido concesiones arrancadas
por una emergencia, y pasada la emergencia se retiraban las concesiones.

\section{La política de esos mismos años}

La reversión no fue solo económica, y el entorno externo no ayudó.

El 24 de febrero de 1996 cazas MiG de la fuerza aérea cubana derribaron dos
avionetas civiles piloteadas por Hermanos al Rescate, una organización del exilio
que buscaba balseros y lanzaba octavillas sobre La Habana, matando a cuatro
hombres. La ubicación del derribo ---dentro o fuera del espacio aéreo cubano--- se
disputó; que las aeronaves eran avionetas civiles desarmadas, no. La consecuencia
política en Washington fue inmediata: la Ley para la Libertad y la Solidaridad
Democrática Cubanas, la Helms-Burton, que estaba estancada, se aprobó en semanas.
Codificó el embargo en un estatuto, de modo que ahora solo puede levantarlo el
Congreso y no un presidente; extendió las sanciones a empresas extranjeras que
traficaran con propiedad confiscada a estadounidenses; y su Título III creó un
derecho de los nacionales norteamericanos a demandar a esas empresas ante
tribunales de Estados Unidos. Todos los presidentes desde Clinton suspendieron el
Título III en tramos de seis meses hasta 2019.

La Helms-Burton importa para la Parte III por una razón que va más allá de su
efecto directo. Convirtió la política norteamericana hacia Cuba de instrumento
ejecutivo, que puede negociarse, en instrumento legislativo, que hay que derogar.
Todo diseño de reforma cubana que dé por supuesto que el embargo terminará cuando
mejoren las relaciones está dando por supuesto algo que la estructura
constitucional norteamericana no entrega con facilidad, y por eso el
capítulo~\ref{ch:premises} trata el embargo como una variable que puede no
moverse.

Dentro de Cuba, esos mismos años produjeron la Ley 88 de 1999, que castigaba con
hasta veinte años la colaboración con la política de Estados Unidos y está
redactada de forma lo bastante amplia como para cubrir casi cualquier forma de
periodismo independiente; y después, en 2002, el Proyecto Varela. Oswaldo Payá y
una red de activistas laicos católicos usaron el Artículo 88 de la constitución de
1976, que permitía a los ciudadanos proponer legislación con 10.000 firmas, para
presentar 11.020 firmas pidiendo un referendo sobre libertad de expresión,
libertad de asociación, una amnistía y elecciones libres. Fue el intento más serio
en cuarenta años de usar las propias reglas del sistema contra él. La respuesta
fue una contramovilización que recogió unos ocho millones de firmas a favor de una
enmienda constitucional que declaraba irrevocable el carácter socialista del
sistema, y que la Asamblea Nacional debidamente aprobó en junio de 2002. La vía
que Payá había encontrado se cerró por enmienda constitucional.

En marzo de 2003, la semana en que Estados Unidos invadió Irak, la seguridad del
Estado cubana detuvo a setenta y cinco periodistas independientes, bibliotecarios,
sindicalistas y organizadores de la oposición, los juzgó en semanas y los condenó
a penas de entre seis y veintiocho años. Se conoce como la Primavera Negra. Ese
mismo mes, tres jóvenes secuestraron una lancha de pasajeros de la bahía de La
Habana intentando llegar a Florida, tomaron rehenes, no hirieron a nadie, se
rindieron cuando la embarcación se quedó sin combustible, y fueron juzgados y
fusilados en nueve días. La posición de Cuba ante la izquierda europea, que había
sobrevivido a Padilla, en buena medida no sobrevivió a esto. Las esposas y madres
de los setenta y cinco empezaron a caminar de blanco a misa cada domingo, y
seguían haciéndolo, y siendo detenidas por ello, veinte años después.

\begin{ledger}{reforma y reversión, 1993--2004}
\emph{Logrado.} Una economía sacada de la caída libre en dieciocho meses. Una
moneda destruida estabilizada sin asistencia externa ---déficit de cerca del 30
por ciento del producto al 2 por ciento, tasa informal de 150 a 25---, que es un
logro técnico genuino y el contraejemplo del capítulo~\ref{ch:crisis}. Dólares,
trabajo por cuenta propia, mercados campesinos e inversión extranjera legalizados.
Turismo reconstruido de la nada hasta 1,8 millones de visitantes. Un sector
privado de unos cientos de miles de personas recreado desde cero, y con él los
primeros cubanos en una generación que no dependían del Estado para su ingreso.

\emph{Costo.} Una dualidad monetaria que volvió ficticia toda la contabilidad
empresarial del país durante veintiséis años, y que con ello volvió la economía
irreformable por cualquiera, incluidos quienes querían reformarla. Desigualdad
desde las más bajas medidas en cualquier parte hasta niveles latinoamericanos en
una década, distribuida a lo largo de una línea racial que el Estado había pasado
treinta años sin querer medir. Enclaves turísticos con entrada legalmente
prohibida a los cubanos. Un programa de reformas retirado en cuanto retirarlo fue
sobrevivible, con el instituto de investigación que lo había analizado
desmantelado como advertencia. Setenta y cinco personas presas con penas de hasta
veintiocho años, y tres fusiladas, en un solo mes.
\end{ledger}
